\section{Distribution Strategy}
\label{sec:distribution}

This section evaluates packaging and distribution models for Timeline Compiler, recommending an architecture that serves human authors, AI agents, and embedding use cases.

\subsection{Distribution Requirements}

Timeline Compiler must be accessible to multiple audiences:

\begin{description}
  \item[CLI Users] — Developers, DevOps engineers, and technical users who invoke rendering from terminal or scripts.
  
  \item[Application Embedders] — Tools that integrate Timeline Compiler as a library (documentation generators, planning tools, dashboards).
  
  \item[AI Agents] — LLM-based agents that need to render timelines as a tool capability.
  
  \item[Non-Technical Users] — Product managers and executives who want live preview while authoring (via IDE integration).
  
  \item[CI/CD Pipelines] — Automated rendering in build processes.
\end{description}

Key requirements:
\begin{itemize}
  \item Zero-friction installation for common platforms
  \item Embeddable as a library
  \item Invocable by AI agents (MCP server or similar)
  \item Deterministic output regardless of invocation method
  \item Self-contained — minimal or no external dependencies at runtime
\end{itemize}

\subsection{Packaging Options}

\subsubsection{Core Library}

\textbf{Concept:} A core rendering library with programmatic API, distributed as a package in the target ecosystem's package manager.

\textbf{Options:}
\begin{description}
  \item[\texttt{@timeline-compiler/core} (npm)] Node.js/TypeScript library. Pros: large ecosystem, easy embedding in JS/TS tools, npm distribution is well-understood. Cons: Node runtime dependency, potentially slow for large renders, binary output (PDF, PNG) may require native dependencies.
  
  \item[Go Module] Single-binary compilation, embeddable via Go's module system. Pros: no runtime dependency, fast, easy cross-compilation. Cons: embedding in non-Go tools requires FFI or subprocess.
  
  \item[Rust Crate] Similar to Go — single binary, high performance. Smaller ecosystem adoption than Go for CLI tools. WASM compilation possible.
  
  \item[Python Package (pip)] Broad adoption in data/ML communities. Cons: runtime dependency, packaging complexity, slower execution.
\end{description}

\textbf{Recommendation:} Implementation language is not decided in this spec, but the core library should expose:
\begin{itemize}
  \item A stable programmatic API (\texttt{render(ir, theme, options) -> output})
  \item Schema validation functions
  \item Theme loading and composition
\end{itemize}

\subsubsection{Command-Line Interface}

\textbf{Concept:} A CLI tool for rendering from terminal.

\begin{lstlisting}[language=bash,caption={CLI usage examples}]
# Basic rendering
timeline render roadmap.yaml -o roadmap.pdf

# With theme
timeline render roadmap.yaml --theme corporate-blue -o roadmap.svg

# Validate only
timeline validate roadmap.yaml

# Watch mode (optional)
timeline watch roadmap.yaml -o roadmap.svg
\end{lstlisting}

\textbf{Distribution:}
\begin{description}
  \item[npm global install] \texttt{npm install -g @timeline-compiler/cli} — Easy for Node.js users, but requires Node runtime.
  
  \item[Standalone Binary] Distributed via GitHub Releases, Homebrew, apt, etc. No runtime dependency. Cross-platform (macOS, Linux, Windows).
  
  \item[Docker Image] \texttt{docker run timeline-compiler render ...} — Isolated environment, good for CI/CD.
  
  \item[npx] \texttt{npx @timeline-compiler/cli render ...} — Zero-install execution for one-off use.
\end{description}

\textbf{Recommendation:} Provide both:
\begin{itemize}
  \item npm package for Node.js ecosystems
  \item Standalone binaries for users without Node.js
\end{itemize}

\subsubsection{VS Code Extension}

\textbf{Concept:} IDE integration providing:
\begin{itemize}
  \item Syntax highlighting for \texttt{.timeline.yaml} files
  \item Live preview panel (re-renders on save or keystroke)
  \item Export commands (PDF, PNG, SVG, PPTX)
  \item Schema validation and error highlighting
  \item Theme selection
\end{itemize}

\textbf{Value:} Dramatically improves authoring experience for non-CLI users. Live preview is critical for iterative refinement.

\textbf{Implementation:} The extension embeds the core library (via WASM, bundled Node module, or subprocess call to CLI).

\textbf{Distribution:} VS Code Marketplace.

\textbf{Recommendation:} High-priority for adoption. Authors should see their timeline update as they type.

\subsubsection{MCP Server (Model Context Protocol)}

\textbf{Concept:} A server exposing Timeline Compiler as a tool for AI agents.

\begin{lstlisting}[language=yaml,caption={MCP tool definition (conceptual)}]
tools:
  - name: render_timeline
    description: "Renders a timeline IR to visual output"
    inputSchema:
      type: object
      properties:
        ir:
          type: object
          description: "Timeline IR conforming to schema"
        theme:
          type: string
          description: "Theme name or inline theme object"
        format:
          type: string
          enum: [svg, png, pdf]
    returns:
      type: string
      description: "Base64-encoded output or file path"
\end{lstlisting}

\textbf{Value:} Agents can render timelines as a native capability. The agent generates IR, invokes the tool, receives visual output.

\textbf{Deployment:} 
\begin{itemize}
  \item Local MCP server (runs alongside the agent environment)
  \item Hosted MCP endpoint (cloud-deployed rendering service)
\end{itemize}

\textbf{Recommendation:} Essential for the agent use case. MCP server should be a first-class distribution target, not an afterthought.

\subsubsection{NPM Package for Embedding}

\textbf{Concept:} The core library packaged for use in JavaScript/TypeScript applications.

\begin{lstlisting}[language=javascript,caption={Embedding in a Node.js application}]
import { render, loadTheme, validate } from '@timeline-compiler/core';

const ir = loadYAML('roadmap.yaml');
const errors = validate(ir);
if (errors.length > 0) throw new ValidationError(errors);

const theme = await loadTheme('corporate-blue');
const svg = await render(ir, theme, { format: 'svg' });

res.setHeader('Content-Type', 'image/svg+xml');
res.send(svg);
\end{lstlisting}

\textbf{Use Cases:}
\begin{itemize}
  \item Documentation sites that render timelines dynamically
  \item Internal tools with timeline visualization
  \item SaaS products embedding Timeline Compiler
\end{itemize}

\textbf{Recommendation:} Core distribution path for the JavaScript ecosystem.

\subsubsection{Standalone Go Binary}

\textbf{Concept:} If the implementation uses Go, the CLI is a single statically-linked binary.

\textbf{Pros:}
\begin{itemize}
  \item No runtime dependencies
  \item Fast startup
  \item Easy cross-compilation (macOS, Linux, Windows, ARM)
  \item Simple distribution (download and run)
\end{itemize}

\textbf{Cons:}
\begin{itemize}
  \item Embedding in non-Go applications requires subprocess or FFI
  \item Community contribution may be lower than TypeScript (smaller Go population among target users)
\end{itemize}

\textbf{Recommendation:} Attractive for CLI distribution. If TypeScript is the core language, a Go wrapper or WASM compilation may achieve similar benefits.

\subsection{Implementation Language Trade-offs}

The distribution strategy depends on implementation language. Key considerations:

\begin{table}[h]
\centering
\caption{Language Trade-offs for Distribution}
\begin{tabularx}{\textwidth}{lXX}
\toprule
\textbf{Factor} & \textbf{TypeScript/Node} & \textbf{Go} \\
\midrule
CLI standalone binary & Requires pkg/nexe or similar & Native \\
npm embedding & Native & Requires WASM or subprocess \\
WASM (browser/edge) & Via esbuild or similar & Native compilation \\
MCP server & Native (Node ecosystem) & Native \\
VS Code extension & Native embedding & Subprocess or WASM \\
Community contributions & Larger pool (frontend devs) & Smaller but growing \\
PDF generation & Needs native dep (puppeteer, etc.) & Native libraries available \\
Startup time & Slower (V8 init) & Fast \\
Binary size & Larger (bundled runtime) & Smaller \\
\bottomrule
\end{tabularx}
\end{table}

\textbf{Recommendation:} This spec does not mandate implementation language. However, the distribution architecture should support:
\begin{enumerate}
  \item A core library embeddable in JavaScript/TypeScript contexts (npm)
  \item A standalone CLI that does not require a runtime install
  \item An MCP server for agent integration
  \item A VS Code extension with live preview
\end{enumerate}

A TypeScript core with WASM or native compilation for CLI may satisfy all requirements. Alternatively, a Go core with a thin Node wrapper for npm distribution is viable.

\subsection{Recommended Architecture}

\begin{verbatim}
                        +-------------------+
                        |   Core Library    |
                        | (render, validate |
                        |   theme, schema)  |
                        +-------------------+
                                 |
         +-----------+-----------+-----------+-----------+
         |           |           |           |           |
    +----+----+ +----+----+ +----+----+ +----+----+ +----+----+
    |   CLI   | |   npm   | |   MCP   | | VS Code | | Docker  |
    |  Binary | | Package | | Server  | |   Ext   | |  Image  |
    +---------+ +---------+ +---------+ +---------+ +---------+
\end{verbatim}

\textbf{Core Library:} Single implementation of IR parsing, validation, theme loading, layout, and rendering. All distribution targets wrap this core.

\textbf{CLI Binary:} Wraps core library. Distributed via:
\begin{itemize}
  \item GitHub Releases (tarballs for each platform)
  \item Homebrew (\texttt{brew install timeline-compiler})
  \item npm global install (if Node.js-based)
  \item curl one-liner for quick install
\end{itemize}

\textbf{npm Package:} The core library published to npm for JavaScript/TypeScript embedding.

\textbf{MCP Server:} Wraps core library, exposes \texttt{render\_timeline} and \texttt{validate\_timeline} tools. Distributed as:
\begin{itemize}
  \item npm package (for local MCP server)
  \item Docker image (for hosted deployment)
\end{itemize}

\textbf{VS Code Extension:} Embeds core library (via bundled WASM or Node module). Published to VS Code Marketplace.

\textbf{Docker Image:} Contains CLI binary. For CI/CD pipelines and isolated execution.

\subsection{Versioning and Compatibility}

\textbf{Semantic Versioning:} All packages follow semver. Breaking changes to IR schema or theme format require major version bump.

\textbf{Schema Versioning:} The IR declares its schema version (\texttt{version: "1.0"}). The compiler reports errors if the IR version is unsupported.

\textbf{Theme Versioning:} Themes declare compatibility with IR/compiler versions.

\textbf{CLI Version Reporting:} \texttt{timeline --version} reports compiler version, supported IR versions, and supported output formats.

\subsection{Distribution Priorities}

For MVP, distribution priority order:

\begin{enumerate}
  \item \textbf{CLI Binary} — Core functionality, scriptable, CI/CD compatible
  \item \textbf{npm Package} — Embedding in JS/TS tools
  \item \textbf{VS Code Extension} — Authoring experience
  \item \textbf{MCP Server} — Agent integration
  \item \textbf{Docker Image} — Isolated execution
\end{enumerate}

Post-MVP:
\begin{itemize}
  \item Homebrew formula
  \item apt/yum packages
  \item GitHub Action for rendering in workflows
  \item Hosted rendering API (optional commercial offering)
\end{itemize}
